\section{Method}
\label{sec:method}

Our main question is: {\bf how does the length and detail of a task instruction affect task performance and format compliance, and does the effect depend on task type and model size?}

\subsection{Instruction tiers}
\label{sec:tiers}

For each task we write five instructions that form a nested ladder.
Let $I_1,\dots,I_5$ denote the \minimal, \concise, \standard, \detailed, and \exhaustive tiers.
We construct $I_k = I_{k-1} \oplus B_k$, where $B_k$ is a new block of text and $\oplus$ is concatenation.
Tier $k$ therefore contains all information in tier $k-1$.
Every tier, including \minimal, states the same core output-format constraint, so compliance is comparable across tiers.
The added blocks are:
\begin{itemize}[leftmargin=*,itemsep=0pt,topsep=0pt]
    \item \minimal: task plus format constraint (18--19 words).
    \item \concise: label definitions or task basics (40--55 words in total).
    \item \standard: numbered guidelines and key decision rules (141--172 words).
    \item \detailed: edge cases, tie-breaking rules, and quality criteria (316--387 words).
    \item \exhaustive: purpose and rationale, a step-by-step procedure, many more rules, a list of common mistakes, a restated format, and a final reminder (820--1{,}018 words).
\end{itemize}
We wrote the \exhaustive tier to look like a realistic ``kitchen-sink'' prompt: long and redundant but internally consistent.
It never contradicts earlier tiers.

\para{Padding control.}
To separate the effect of more tokens from the effect of more content, we build a \padded condition: the \standard text followed by task-irrelevant office boilerplate (a narrative about a documentation program) truncated so that the total word count matches \exhaustive.
The filler contains no task instructions and no label information.
If long instructions hurt because of token count alone, \padded should behave like \exhaustive; if content drives the effect, \padded should behave like \standard.
\Tabref{tab:tiers} lists word and token counts for all six conditions.

\begin{table}[t]
\centering
\caption{Instruction length per condition. Words count the instruction only; tokens count the full prompt including the task input (mean over items). The tiers are nested; \padded is word-matched to \exhaustive and has 7--12\% fewer tokens.}
\label{tab:tiers}
\small
\begin{tabular}{@{}lcccccccc@{}}
\toprule
& \multicolumn{4}{c}{Instruction words} & \multicolumn{4}{c}{Prompt tokens (incl.\ input)} \\
\cmidrule(lr){2-5} \cmidrule(lr){6-9}
Condition & \textsc{Cls} & \textsc{Ext} & \textsc{Gen} & \textsc{Rsn} & \textsc{Cls} & \textsc{Ext} & \textsc{Gen} & \textsc{Rsn} \\
\midrule
\minimal & 19 & 18 & 19 & 18 & 112 & 230 & 621 & 117 \\
\concise & 55 & 50 & 50 & 40 & 172 & 273 & 658 & 142 \\
\standard & 163 & 172 & 165 & 141 & 323 & 439 & 809 & 270 \\
\detailed & 387 & 360 & 343 & 316 & 605 & 672 & 1032 & 486 \\
\exhaustive & 1018 & 866 & 897 & 820 & 1461 & 1325 & 1723 & 1091 \\
\midrule
\padded & 1018 & 866 & 897 & 820 & 1291 & 1225 & 1638 & 1038 \\
\bottomrule
\end{tabular}
\end{table}

\subsection{Tasks and data}

We use four task types, each with one public dataset.
For each task we draw a fixed random sample of $n=252$ items (seed 42) and use the same items in every condition, so all comparisons are paired.
\begin{itemize}[leftmargin=*,itemsep=0pt,topsep=0pt]
    \item {\bf Classification}: \agnews~\citep{zhang2015character} test set, stratified to 63 items per class. Format: answer with only one of World, Sports, Business, Sci/Tech. Metric: accuracy (chance is 25\%).
    \item {\bf Extraction}: \squad v1.1~\citep{rajpurkar2016squad} validation set, one question per passage. Format: only a span copied verbatim from the passage. Metric: token F1.
    \item {\bf Generation}: \cnndm 3.0.0~\citep{hermann2015teaching,see2017get} test set, articles truncated to 500 words. Format: 2--4 sentences, at most 80 words, summary only. Metric: ROUGE-L F1~\citep{lin2004rouge} against the reference highlights.
    \item {\bf Reasoning}: \gsm~\citep{cobbe2021training} test set. Format: the last line is \texttt{Final answer: <number>}. Metric: exact-match accuracy.
\end{itemize}

\subsection{Models and decoding}

No API access was available, so we run five open-weight instruction-tuned models locally: \qwen-\{0.5B, 1.5B, 3B, 7B, 14B\}-Instruct~\citep{qwen2025qwen25}, all in bf16, served with vLLM~\citep{kwon2023efficient}.
We use greedy decoding (temperature 0, seed 42), the default chat template, and place the instruction in the user turn before the task input.
The maximum number of new tokens is 16 (classification), 64 (extraction), 200 (generation), and 768 (reasoning), with a 4{,}096-token context.
Each model processes $6 \times 4 \times 252 = 6{,}048$ prompts.
As a robustness check, we rerun 0.5B and 1.5B with the instruction in the system message.
We also planned to run Llama-3.1-8B-Instruct and \qwen-32B-Instruct, but other jobs occupied the shared GPUs for the rest of the session and these runs could not start.

\subsection{Scoring}

We separate \emph{what} the model answers from \emph{how} it formats the answer.

\para{Task score (lenient parsing).}
A formatting slip should not count as a wrong answer.
For classification we take the first label or synonym mentioned; for extraction, the first line with ``Answer:'' prefixes and quotes removed; for reasoning, the number after ``Final answer,'' falling back to the last number in the output.

\para{Compliance (strict parsing).}
Classification output must be exactly a label.
Extraction output must be one line, a verbatim substring of the passage, at most 30 words.
Generation output must have 2--4 sentences, at most 80 words, and no bullets, markdown, or preamble such as ``Here is\ldots''.
Reasoning output must end with a line that matches \texttt{Final answer: <number>} exactly.

\subsection{Hypotheses and statistics}
\label{sec:hypotheses}

We pre-registered the following hypotheses:
{\bf H1}: classification and extraction peak at \standard and drop 3--7 points at \exhaustive;
{\bf H2}: generation improves monotonically up to \detailed;
{\bf H3}: reasoning does not change beyond \concise;
{\bf H4}: 7B degrades more from \exhaustive than larger models;
{\bf H5}: compliance peaks at \standard for every task;
{\bf H6}: comparing \padded with \standard and \exhaustive separates token dilution from content.

We compute 95\% bootstrap confidence intervals over items (5{,}000 resamples) for each model$\times$task$\times$condition cell.
For paired differences we report bootstrap CIs and two-sided sign-flip permutation tests.
We apply Holm correction within the primary family of 20 \standard\versus\exhaustive tests (5 models $\times$ 4 tasks) and within the 40 padding-control tests.
For H2 we fit an item-level linear slope over tiers \minimal to \detailed.
For H3 we use Cochran's Q over \concise to \exhaustive, plus an equivalence check (90\% CI within $\pm$3 points).
For H4 we compute an item-level paired difference-in-differences between 7B and 14B, and the Spearman correlation between the \standard$\to$\exhaustive drop and log model size.
